\label{der:4-2}
对$[A\mid\vb]$仅作行变换，先把系数块化为阶梯形。若某个系数全零的行右端为$c\ne0$，对应$0=c$，无解；该右端位置正是增广列新增的主元，所以增广秩比系数秩大。若没有这种行，全部非零系数行均有主元，任给非主元变量，都可从最后一个主元行向上唯一回代，所以有解，增广列也不新增主元。相容性与两秩相等因而等价。

设系数秩$r$，主元变量$r$个，自由变量$n-r$个。齐次时令第$j$个自由变量为1、其他自由变量为0，回代得到$\eta_j$。其自由分量组成单位阵，所以这些向量无关。任意齐次解的自由分量为$c_1,\ldots,c_{n-r}$，减去$\sum_jc_j\eta_j$后自由分量全零，回代使主元分量也全零，因此它们生成全部齐次解，构成基础解系。

相容时选任一特解$\vx_p$。若$\vx$也是解，则$A(\vx-\vx_p)=\vb-\vb=0$，因此$\vx-\vx_p$是上面基础解系的组合。反之，$A(\vx_p+\sum c_j\eta_j)=\vb$，所以式\eqref{eq:general}恰好给全部解。$r=n$时没有自由变量，只有特解；$r<n$时在实数、复数域有无穷多组参数，得到无穷多解。
